\begin{apartats}

\apartat[2,5]{1,5}
Determina $a$, $b$ i $c$ perquè la funció
\[
f(x)=x^3+ax^2+bx+c
\]
passi pel punt $(0,2)$, tingui un punt d'inflexió en $x=1$ i un extrem relatiu en $x=3$.

\begin{solucio}
Derivades: $f'(x)=3x^2+2ax+b$ i $f''(x)=6x+2a$.\\
\textbf{Passa per $(0,2)$}: $f(0)=c=2$, d'on $\boxed{c=2}$.\\
\textbf{Inflexió en $x=1$}: $f''(1)=6+2a=0$, d'on $\boxed{a=-3}$.\\
\textbf{Extrem en $x=3$}: $f'(3)=27+2(-3)(3)+b=27-18+b=0$, d'on $\boxed{b=-9}$.\\
La funció és $f(x)=x^3-3x^2-9x+2$. Comprovació: $f''(1)=0$ i $f'(3)=27-18-9=0$.
\end{solucio}

\begin{nomesllarg}
\apartat{1}
Amb aquests valors, digues si l'extrem de $x=3$ és un màxim o un mínim, i estudia la curvatura
de $f$.

\begin{solucio}
$f''(x)=6x-6$, i $f''(3)=12>0$: en $x=3$ hi ha un \textbf{mínim relatiu}.\\
$f''<0$ a $(-\infty,1)$: \textbf{còncava}. \quad $f''>0$ a $(1,+\infty)$: \textbf{convexa}.\\
La curvatura canvia en $x=1$, que és el punt d'inflexió demanat.
\end{solucio}
\end{nomesllarg}

\end{apartats}
